\documentclass[11pt]{article}
\usepackage[margin=.85in]{geometry}
\usepackage[T1]{fontenc}
\usepackage{lmodern,microtype,amsmath,amssymb,graphicx,booktabs,tabularx}
\usepackage{xcolor}
\usepackage[colorlinks=true,linkcolor=blue!40!black,urlcolor=blue!50!black]{hyperref}
\setlength{\parindent}{0pt}
\setlength{\parskip}{.4em}
\setlength{\emergencystretch}{1.5em}
\graphicspath{{images/}}
\title{Observed cluster tests of the conservative stock sector}
\author{Research calculation for the Entropic Scalar EFT project}
\date{23 October 2026}
\begin{document}
\maketitle
\begin{abstract}
Public radial profiles for twelve X-COP clusters, five matching weak-lensing mass estimates and a released Bullet Cluster lens reconstruction provide the first resolved observational targets for the conservative stock dynamics. The existing surviving cohort allocation is held fixed. A restricted model in which stock follows the present baryons is insufficient across these targets. A2142 is close to its mass budget; A1795 requires substantial relative segregation if its independent lensing mass is adopted. The Bullet reconstruction supplies simultaneous aperture-mass, gas-separation and galaxy-alignment targets. These calculations constrain the required outcome of transport. They do not demonstrate that the proposed interaction produces it.
\end{abstract}

\section{The tested statement}
The preceding encounter calculation evolved gravitating stock, stellar tracers and baryonic gas with reciprocal impulses and a finite-work contact. Its initial profiles and coupling parameters were specified. Here the gravitating profiles come from observations, and their associated assumptions remain explicit.

We retain the project's saturated surviving allocation
\begin{equation}
k=5.1760752079627395,\qquad f_{b,\mathrm{coext}}=\frac{1}{1+k}=0.1619151.
\end{equation}
It is a stock-to-\emph{original-cohort}-baryon ratio. For a restricted comparison, impose the stronger assumption that stock and present baryons occupy each aperture in that ratio:
\begin{equation}
M_{\mathrm{coext}}(<r)=(1+k)M_b(<r),\qquad M_b=M_g+M_\star.
\label{eq:shortcut}
\end{equation}
Equation~\eqref{eq:shortcut} is a testable shortcut, not a necessary consequence of the theory. A gravitating target instead requires
\begin{equation}
M_{c,\mathrm{req}}=M_{\mathrm{grav}}-M_b,\qquad
\eta(r)=\frac{kM_b}{M_{\mathrm{grav}}-M_b}.
\label{eq:eta}
\end{equation}
For a single conserved cohort with aperture retentions $r_b$ and $r_c$, $\eta=r_b/r_c$. In a mixed aperture it is an effective inventory diagnostic: imported material must be tracked separately. Values above one can reflect preferential stock loss or imported baryons. Neither sign of $\eta-1$ alone establishes a violation of the whole-cohort inventory.

The normalization is unchanged between clusters. The full-release convention is itself a conditional input; it is not evidence that each observed cluster has reached the same release state. No per-cluster coupling, stock amplitude or acceleration scale is fitted here.

\clearpage
\section{Resolved X-COP profiles}
We downloaded the complete primary X-COP archive and read its FITS tables for all twelve clusters~\cite{xcop,ettori,eckert}. Seven have released stellar profiles; the other five receive gas-only lower-bound diagnostics. No stellar mass is invented for those five. The comparison uses the released hydrostatic masses without a correction that enforces the cosmic baryon fraction.

\begin{center}
\includegraphics[width=.98\linewidth]{xcop_radial.png}
\end{center}
{\small Solid blue: equation~\eqref{eq:shortcut} divided by the NFW hydrostatic mass. Dashed orange: the same baryons divided by the forward hydrostatic reconstruction. Blue shading spans zero to the measured projected stellar mass, for the measured stellar component only; it is not a confidence band. Dotted blue doubles the adopted deprojected stellar contribution. This is a sensitivity scenario for stellar-population and diffuse-light uncertainty, not a hard upper bound.}

All principal radial tests start at $0.1R_{500}$ and stop within the available gas, mass and stellar profiles, with a maximum of $1.5R_{500}$. The measured shapes generally require more centrally concentrated stock than a local rescaling of the gas and stars. Agreement near one outer aperture does not reproduce the inner profile. Correlated radial samples are not summed into a fictitious independent-bin $\chi^2$.

\clearpage
\section{Mass inventory and reconstruction sensitivity}
The stellar FITS profiles are projected cumulative masses. Ghizzardi et al. apply a factor $0.75$ at $R_{500}$ to obtain a spherical mass and exclude intracluster light~\cite{ghizzardi}. We preserve that distinction. Our radial extension uses the geometric sphere-to-cylinder ratio of a spherical generalized NFW distribution with inner slope $1.64$, choosing its scale once so that $D(R_{500})=0.75$. This gives $R_{500}/r_s=1.28214$. Multiplying the measured profile by this radius-dependent factor is an approximation, not a recovered three-dimensional stellar distribution. The reported geometric and doubled-star sensitivities bracket alternatives for this measured component.

\begin{center}
\begin{tabular}{lrrr}
\toprule
Cluster & $M_{\mathrm{coext}}/M_{\mathrm{HSE}}$ & $\eta_{\mathrm{HSE}}$ & Ratio with $2M_\star$\\
\midrule
A1795 & 0.903 & 0.887 & 0.943\\
A2029 & 0.912 & 0.897 & 0.958\\
A2142 & 1.023 & 1.027 & 1.071\\
A2319 & 1.213 & 1.265 & 1.257\\
A644  & 0.851 & 0.827 & 0.897\\
A85   & 0.952 & 0.943 & 0.976\\
ZW1215& 0.679 & 0.640 & 0.705\\
\bottomrule
\end{tabular}
\end{center}
All entries are evaluated at the released hydrostatic $R_{500}$. The source tables, interpolation results and full precision are saved alongside this note.

The nominal NFW reconstruction requires positive cumulative stock and positive stock in each of fifteen finite radial shells for every cluster tested. This is a conditional consistency result. The NFW mass shape was supplied by the observational reconstruction, and positive shell masses do not prove a nonnegative distribution function, stable equilibrium or formation history.

The less constrained forward reconstruction for A644 gives four negative residual shell masses over the tested range. We retain this failure. It prevents a blanket claim that arbitrary released hydrostatic profiles admit a positive stock density, and points to sensitivity to the mass reconstruction and hydrostatic assumptions.

\begin{center}
\includegraphics[width=.90\linewidth]{retention_profiles.png}
\end{center}
{\small The required relative retention changes with radius. These curves are inferred targets; selecting them as an interaction's equilibrium would not constitute a prediction.}

\clearpage
\section{Five independent weak-lensing comparisons}
We matched A1795, A2029, A2142, A85 and ZW1215 to the public MENeaCS results~\cite{herbonnet}. The lensing masses do not use the X-ray hydrostatic equilibrium equation. They retain their authors' NFW concentration prescription and simulation calibration. The aperture and NFW methods use the same shear observations and are not counted as independent datasets.

For a mass prediction under equation~\eqref{eq:shortcut}, solve for its own overdensity radius:
\begin{equation}
(1+k)M_b(<R_{500}^{\mathrm{coext}})
=\frac{4\pi}{3}\,500\rho_{\mathrm{crit}}(z)
\bigl(R_{500}^{\mathrm{coext}}\bigr)^3.
\end{equation}
Thus the prediction is not evaluated at a conveniently selected hydrostatic radius. Separately, equation~\eqref{eq:eta} is evaluated at the radius corresponding to each measured lensing mass. Its one-sigma endpoint diagnostics also change the aperture radius.

\begin{center}
\begin{tabular}{lrrrr}
\toprule
 & \multicolumn{3}{c}{$M_{500}\ [10^{14}M_\odot]$} & $\eta$ at lensing $R_{500}$\\
Cluster & Coextensive & Lensing NFW & Lensing aperture & NFW / aperture\\
\midrule
A1795 & 3.96--4.23 & $9.3\pm2.2$ & $9.2^{+2.3}_{-2.3}$ & 0.518 / 0.523\\
A2029 & 7.46--8.07 & $12.1\pm2.5$ & $10.8^{+1.9}_{-2.3}$ & 0.704 / 0.761\\
A2142 & 9.29--10.02 & $9.7\pm2.3$ & $11.3^{+2.1}_{-2.2}$ & 0.969 / 0.867\\
A85 & 5.20--5.43 & $5.7\pm2.2$ & $7.4^{+2.3}_{-2.3}$ & 0.937 / 0.779\\
ZW1215 & 4.19--4.46 & $3.5\pm2.2$ & $7.0^{+2.4}_{-1.4}$ & 1.143 / 0.686\\
\bottomrule
\end{tabular}
\end{center}
The coextensive range changes the adopted stellar mass from one to two times its nominal value. It is not a statistical interval.

\begin{center}
\includegraphics[width=.95\linewidth]{independent_lensing.png}
\end{center}
A1795 is the most useful stress test in this overlap sample: its two lensing estimators agree closely and require roughly half as much retained baryonic mass per unit retained stock as a coextensive cohort. Its nominal mass difference is $2.43$ times the quoted NFW lensing error alone ($2.30$ with doubled stars). These ratios are descriptive residuals, not a combined exclusion significance. A2029 also requires relative segregation. ZW1215 is sensitive to the lensing estimator, so its hydrostatic deficit cannot be promoted to an equally robust lensing discrepancy.

\clearpage
\section{The Bullet Cluster as a spatial target}
We retrieved the current public CANUCS Bullet v1 best-fit convergence map and its fixed gas component~\cite{bullet,canu}. Their documentation sets $D_{ls}/D_s=1$, giving
\begin{equation}
\Sigma_{\mathrm{crit}}=\frac{c^2}{4\pi G D_l}
=1.8268761\times10^9\ M_\odot\,\mathrm{kpc}^{-2},
\qquad 1''=4.41302\ \mathrm{kpc}
\end{equation}
for the release cosmology at $z_l=0.296$. We integrate matched circular apertures after registering the FITS world coordinates. The total reconstruction already contains the supplied gas. Subtracting gas leaves a non-gas component that still includes stars; galaxy lens halos are not pure stellar mass.

\begin{center}
\includegraphics[width=\linewidth]{bullet_observed_maps.png}
\end{center}
\begin{center}
\begin{tabular}{lrrrr}
\toprule
Centre & Radius (kpc) & $M_{\rm total}/10^{13}M_\odot$ & $M_g/10^{13}M_\odot$ & $M_g/M_{\rm total}$\\
\midrule
BCG1 & 60 & 2.315 & 0.163 & 0.070\\
BCG2 & 60 & 2.272 & 0.206 & 0.091\\
BCG3 (Bullet) & 60 & 2.151 & 0.099 & 0.046\\
BCG1 & 150 & 11.281 & 1.006 & 0.089\\
BCG2 & 150 & 11.607 & 1.259 & 0.108\\
BCG3 (Bullet) & 150 & 8.511 & 0.618 & 0.073\\
\bottomrule
\end{tabular}
\end{center}
The released subcluster halo centre is $3.92$ kpc from BCG3, reproducing the published roughly $4$ kpc alignment. The supplied gas-density map has its local peak about $169$ kpc from BCG3; this is a native-pixel diagnostic with an $8.68$ kpc gas pixel, not an X-ray brightness-peak measurement or a fitted offset posterior. Aperture integration changes by at most $0.214\%$ when pixel sampling is coarsened. Observational and model uncertainties are larger.

These are measurements of an existing lens reconstruction. No simulated map has been optimized against it. The previous encounter's $24$--$30$ kpc force softening, much broader map smoothing, different mass ratio and specified Plummer profiles cannot establish agreement at the observed few-kiloparsec alignment.

\clearpage
\section{What the effective dynamics must now accomplish}
Three complementary targets emerge. A2142 provides a useful normalization and equilibrium benchmark. A1795 demands a quantitative assembly and transport explanation while retaining the observed gas profile. The Bullet demands separated gas and gravitating concentrations together with close halo--galaxy alignment. A model should meet these requirements using the same constitutive coupling, with environmental state changes computed from the work and depletion variables.

The useful improvement is a set of numerical obligations. A nonnegative stock inventory can be written for the nominal NFW targets, and the previous conservative experiment supplies a candidate transport mechanism. Connecting them requires a selected phase-space distribution and coupling law that survive an isolated-cluster control and then the merger. The strongest previously tested contact changed the isolated halo appreciably, so increasing its rate is not an established solution to A1795's retention target. Depletion can suppress the interaction, but assigning a favorable depleted state to each cluster would introduce a hidden fit unless its history is computed.

A next forward calculation should fix gas and stellar inputs, specify the initial cohort and formation history, and derive the stock evolution under the common potential. Lensing profiles, projected gas morphology and galaxy offsets then provide the comparison. Fitting the observed total mass first and calling the recovered stock a prediction would not perform this test. A contact-rate bound also requires this forward calculation; a particle-dark-matter self-interaction limit cannot simply be relabeled as the rate of the present stock--baryon process.

\section{Scope and reproducibility}
The analysis adopts the observational $H_0=70\ \mathrm{km\,s^{-1}\,Mpc^{-1}}$, $\Omega_m=0.3$ distance convention. Small redshift differences between the two catalogues are retained in the source record; the largest distance mismatch is about $2\%$ for ZW1215. We use the X-COP redshift for radius conversion and do not claim a recalibration of the original shear-source efficiencies. Stellar deprojection, missing diffuse light, hydrostatic departures, triaxiality, correlated profile errors and lens-model systematics remain relevant. No joint likelihood or new theorem of microphysical uniqueness is claimed.

The project directory contains the pre-calculation protocol; the original X-COP archive; 115 expanded FITS tables; the primary lensing paper and transcribed five-row table; the Bullet maps, catalogues and parameter files; acquisition URLs and checksums; twelve radial result tables; and separate radial and Bullet summaries. An oversized whole-folder Bullet download was stopped after the required small input files were recovered and individually CRC-checked. Its partial archive is preserved, and the complete convergence map was retrieved directly with a recorded checksum.

The calculation scripts are \texttt{radial\_analysis.py} and \texttt{bullet\_analysis.py}. \texttt{validate.py} checks independent critical-density normalization, overdensity roots, cohort normalization, stellar projection convention, world-coordinate round trips, map subtraction, aperture convergence and source hashes. All fifteen numerical/data-handling checks pass. The A644 forward-profile failure is explicitly retained as an analysis result. These checks validate the reduction, not the theory. The four publication papers and preceding simulation outputs are unchanged.

\begin{thebibliography}{9}\small
\bibitem{xcop} X-COP collaboration. \href{https://dominiqueeckert.wixsite.com/xcop/data}{Public data release}, archive linked 22 July 2025; retrieved 23 October 2026.
\bibitem{ettori} S. Ettori et al. \emph{Hydrostatic mass profiles in X-COP galaxy clusters}. A\&A 621, A39 (2019). \href{https://arxiv.org/abs/1805.00035}{arXiv:1805.00035}.
\bibitem{eckert} D. Eckert et al. \emph{Non-thermal pressure support in X-COP galaxy clusters}. A\&A 621, A40 (2019). \href{https://arxiv.org/abs/1805.00034}{arXiv:1805.00034}.
\bibitem{ghizzardi} S. Ghizzardi et al. \emph{Iron in X-COP}. A\&A 646, A92 (2021), Sec. 4.1. \href{https://arxiv.org/abs/2007.01084}{arXiv:2007.01084}.
\bibitem{herbonnet} R. Herbonnet et al. \emph{CCCP and MENeaCS: (updated) weak-lensing masses for 100 galaxy clusters}. MNRAS 497, 4684 (2020), Table A2. \href{https://arxiv.org/abs/1912.04414}{arXiv:1912.04414}.
\bibitem{bullet} G. Rihtar\v{s}i\v{c} et al. \emph{Mapping dark matter in the Bullet Cluster using JWST imaging and spectroscopy} (2026). \href{https://arxiv.org/html/2601.22245v2}{arXiv:2601.22245v2}.
\bibitem{canu} CANUCS. \href{https://niriss.github.io/lensing.html}{Public Bullet lens-model release}, v1 README dated 17 March 2026; retrieved 23 October 2026.
\end{thebibliography}
\end{document}
